\section{批量大小与学习率的 Scaling}

\subsection{Critical Batch Size}

批量大小（batch size）是一个同时涉及系统效率和优化效率的超参数。

\begin{figure}[H]
\centering
\includegraphics[width=0.95\textwidth]{slides-images/slide-024.jpg}
\caption{批量大小的收益递减效应\protect\footnotemark}
\end{figure}
\footnotetext{Slides 第25页。}

批量大小的关键行为分为两个阶段：

\begin{enumerate}
    \item \textbf{完美缩放阶段}：当 batch size $< $ noise scale（梯度噪声尺度），增大 batch size 几乎等效于增加梯度步数——这是最理想的状态，因为并行化带来了真正的加速
    \item \textbf{收益递减阶段}：当 batch size $>$ noise scale，额外的样本不再有效减少噪声，性能改善被优化景观的曲率限制
\end{enumerate}

\textbf{Critical Batch Size} 就是这两个阶段之间的转折点。

\begin{figure}[H]
\centering
\includegraphics[width=0.95\textwidth]{slides-images/slide-025.jpg}
\caption{Critical Batch Size 随目标损失的变化：损失越低，可用的 batch size 越大\protect\footnotemark}
\end{figure}
\footnotetext{Slides 第26页。}

一个非常有趣的发现是：\textbf{目标损失越低，Critical Batch Size 越大}。直觉是：当你追求更低的损失时，模型处于更``精细''的优化区域，梯度噪声的影响更大，因此需要更大的 batch size 来去噪。

\begin{knowledgebox}{LLaMA 3 的 Batch Size 策略}
LLaMA 3 训练报告中实际采用了``逐步增大 batch size''的策略。这正是基于 Critical Batch Size 随目标损失变化的洞察：训练初期（损失高）用较小 batch size，训练后期（损失低）增大 batch size，兼顾优化效率和系统效率。
\end{knowledgebox}

\subsection{Batch Size 与 Compute 的关系}

\begin{figure}[H]
\centering
\includegraphics[width=0.95\textwidth]{slides-images/slide-026.jpg}
\caption{Kaplan：随计算量增加，最优 batch size 增大但步数可保持不变\protect\footnotemark}
\end{figure}
\footnotetext{Slides 第27页。}

Kaplan 的分析表明：随着计算预算增大，最优 batch size 也增大，但\textbf{总训练步数可以基本保持不变}。这对数据并行系统来说是好消息——意味着可以通过增加 GPU 数量（增大 batch size）而不是增加训练时间来利用更多计算资源。

\subsection{学习率 Scaling 与 MuP}

学习率是另一个与模型规模紧密相关的超参数。

\begin{figure}[H]
\centering
\includegraphics[width=0.95\textwidth]{slides-images/slide-027.jpg}
\caption{标准做法 vs. MuP：学习率随模型宽度的变化\protect\footnotemark}
\end{figure}
\footnotetext{Slides 第28页。}

\textbf{标准做法}下，最优学习率随模型宽度增大而\textbf{递减}。经验法则是学习率 $\propto 1/\text{width}$。这意味着每次改变模型大小都需要重新搜索最优学习率。

\textbf{MuP（Maximal Update Parameterization）} 提供了一种更优雅的解决方案：通过重新参数化模型——调整不同层的初始化方差和学习率缩放，使得\textbf{最优学习率在不同宽度下保持不变}。

\begin{importantbox}{MuP 的核心思想}
MuP 通过以下方式重新参数化模型：
\begin{itemize}
    \item 根据模型宽度缩放每层的初始化方差
    \item 根据模型宽度缩放每层的学习率
    \item 在前向传播中乘以宽度相关的缩放因子
\end{itemize}
效果是：在最小模型上调好的学习率，\textbf{直接迁移到最大模型}，无需重新搜索。
\end{importantbox}

\begin{knowledgebox}{MuP 的行业采用}
Meta 在发布 LLaMA 4 时声称使用了一种叫做 ``MetaP'' 的类似技术。越来越多的实验室认识到，如果不得不依赖预测最优学习率的 Scaling Law 拟合，那么拟合误差可能导致次优结果。通过 MuP 这样的重新参数化，可以从根本上避免这个问题。
\end{knowledgebox}

\subsection{本章小结}

Batch size 和学习率是两个与模型规模高度耦合的超参数。Critical Batch Size 的概念提供了一个理解``何时可以安全增大 batch size''的框架。MuP 等重新参数化方法则从根本上解决了学习率跨尺度迁移的问题。在实践中，许多前沿实验室已经采用了这些技术来优化训练流程。

%% ========== Part 6: 注意事项 ==========

